% ===== 原文 pp.156–171（PDF pp.163–178） =====
% [原文 p.156]
得到的带关系式路代数 $(K\vec{\Delta};\{\rho_{kl}\})$ 与多项式环 $K[x_1,\dots,x_n]$ 同构。注意关系式对应于变量的可交换性。

箭图 $\vec{\Delta}$ 的带关系式 $K$ 上路代数是基本代数。其逆也已知成立。

\begin{theorem}{6.31}
（参看 [2]、[5]。）设 $K=\C$（或 $K$ 为代数闭域）。对 $K$ 上基本代数 $A$，存在某箭图 $\vec{\Delta}$ 的路代数 $K\vec{\Delta}$ 与其关系式 $\rho$，使得 $A\simeq(K\vec{\Delta},\rho)$。
\end{theorem}

\subsection{带关系式箭图的表示论}

\begin{definition}{6.32}
对箭图 $\vec{\Delta}$ 与其路代数 $K\vec{\Delta}$ 的关系式 $\rho$，$(\vec{\Delta},\rho)$ 的表示 $V:=(\bigoplus_{i\in\Delta_0}V_i,\bigoplus_{\alpha\in\Delta_1}B_\alpha)$，是指 $\vec{\Delta}$ 的表示（定义 6.20），且对每个关系式 $\sum_{k=1}^{r}c_kw_k$（$w_k=(i|\alpha_{i_1}\cdots\alpha_{i_k}|j)$）满足
\[
\sum_{k=1}^{r}c_k\bigl(B_{\alpha_{i_k}}\circ\cdots\circ B_{\alpha_{i_1}}\bigr)=0
\]
者。
\end{definition}

% [Tutor Note: 底本 p.156 定義 6.32 第 3、4 行将路的箭与合成中的首项印作 $\alpha_1$；据作者正誤表（2014-05-20 版）订正为 $\alpha_{i_1}$，此处按订正后的内容翻译。原印 $w_k=(i|\alpha_1\cdots\alpha_{i_k}|j)$、$B_{\alpha_{i_k}}\circ\cdots\circ B_{\alpha_1}$ 中的 $\alpha_1$ 与下标 $i_k$ 不相配（$w_k$ 的第 $1$ 个箭被记为 $\alpha_{i_1}$ 才与末项 $\alpha_{i_k}$ 成同一套指标）。]

$(\vec{\Delta},\rho)$ 的表示所成的范畴 $\mathrm{Rep}(\vec{\Delta},\rho)$ 成为阿贝尔范畴；与没有关系式的情形完全同样地，以下定理成立。

\begin{theorem}{6.33}
令 $A:=(K\vec{\Delta},\rho)$，则阿贝尔范畴的等价
\[
\modA\simeq\mathrm{Rep}(\vec{\Delta},\rho)
\]
成立。
\end{theorem}

\section{导出范畴 $D^{b}(\mathrm{mod}\text{-}A)$ 的生成系及其变异}

对附属于有限型有序箭图的带关系式路代数 $A$，讨论阿贝尔范畴 $\modA$ 与导出范畴 $\Db(\modA)$ 的结构。

\subsection{$\modA$ 的单对象}

\begin{definition}{6.34}
在阿贝尔范畴 $\cC$ 中，对象 $M\in\cC$ 是单对象（simple），是指不存在既非零对象、也非与 $M$ 自身同构的对象 $N\in\cC$ 到 $M$ 的单射 $N\to M$。
\end{definition}

设 $A$ 为某有限型有序箭图 $\vec{\Delta}=(\Delta_0,\Delta_1)$、$\Delta_0=\{1,\dots,n\}$ 的某带关系式路代数。取 $\Delta_0$ 的元的标号，使得对所有 $\alpha\in\Delta_1$ 都有
\[
t(\alpha)<s(\alpha)
\]
成立。对每个 $i\in\Delta_0$，$A$ 中存在本原幂等元 $e_i:=(i|i)$，而它们给出本原正交幂等元分解 $1_A=e_1+\cdots+e_n$。

% [原文 p.157]
把 $\modA$ 的对象与对应的带关系式箭图 $\vec{\Delta}$ 的表示所成阿贝尔范畴 $\mathrm{Rep}(\vec{\Delta})$ 的对象视为同一，把 $M\in\modA$ 作为 $K$ 上向量空间表为
\[
M=\bigoplus_{i\in\Delta_0}M_i,\qquad M_i=M\cdot e_i.
\]
进一步，令
\[
F^{k}(M):=\bigoplus_{i\le k}M_i,
\]
则作为 $K$ 向量空间有 $F^{k}(M)\subset F^{k+1}(M)$，并得到对应的单射列
\[
0=F^{0}(M)\to F^{1}(M)\to\cdots\to F^{n-1}(M)\to F^{n}(M)=M,
\]
但实际上这些单射都是右 $A$ 模的同态映射。也就是说，$F^{k-1}(M)$ 是 $F^{k}(M)$ 的部分 $A$ 模这一点，由对 $A$ 的基 $\alpha$ 必有 $t(\alpha)<s(\alpha)$ 成立而得到。

现在固定 $k\in\Delta_0$，则得到交换图表
\[\begin{tikzcd}[row sep=2.6em, column sep=2.6em]
  0 \arrow[r] & F^{k-1}(M) \arrow[r,"\id_{F^{k-1}(M)}"] \arrow[d,"0"'] & F^{k-1}(M) \arrow[r,"0"] \arrow[d,"B"'] & 0 \arrow[r] \arrow[d,"0"'] & 0 \\
  0 \arrow[r] & 0 \arrow[r,"0"] & M_k \arrow[r,"\id_{M_k}"] & M_k \arrow[r,"0"] & 0 .
\end{tikzcd}\]
这里 $B\colon M_k\to F^{k-1}(M)$ 是把 $\xi\in M_k$ 对应到 $\bigoplus_{s(\alpha)=k}\xi\cdot\alpha\in F^{k-1}(M)$ 的映射。把该图表的上段与下段合起来，看作 $F^{k-1}(M)=F^{k-1}(M)\oplus 0$、$F^{k}(M)=F^{k-1}(M)\oplus M_k$、$M_k=0\oplus M_k$，便得到右 $A$ 模的短正合列
\begin{equation}
0\to F^{k-1}(M)\to F^{k}(M)\to M_k\to 0. \tag{6.3}
\end{equation}
这里，令 $M_k$ 作为 $K$ 向量空间的维数为 $d_k$，则 $M_k\simeq K^{\oplus d_k}$。

\begin{definition}{6.35}
对每个 $k\in\Delta_0$，如下定义 $\modA$ 的单对象 $S(k)\in\modA$：
\[
S(k)\cdot e_i=
\begin{cases}
K\ (\text{$1$ 维向量空间}), & i=k,\\
0, & i\neq k.
\end{cases}
\]
\end{definition}

也就是说，在 (6.3) 中 $M_k\simeq(S(k))^{\oplus d_k}$。

$S(1),\dots,S(n)$ 是单对象这一点是显然的。

\begin{lemma}{6.36}
$\modA$ 的单对象仅限于 $S(1),\dots,S(n)$。
\end{lemma}

\begin{proof}
对某对象 $M\in\modA$，设 $M_k\neq 0$ 的最小 $k$ 为 $k_0$，则 $M$ 含有部分 $A$ 模 $M_{k_0}$；令 $K$ 向量空间 $M_{k_0}$ 的维数为 $d_{k_0}$，则 $M_{k_0}\simeq(S(k_0))^{\oplus d_{k_0}}$。也就是说，当 $M$ 不是零对象时，若 $M$ 是

% [原文 p.158]
$S(1),\dots,S(n)$ 中某一个以外者，则必存在既非零对象也非 $M$ 自身的部分 $A$ 模。因此可知 $S(1),\dots,S(n)$ 以外的单对象不存在。
\end{proof}

\subsection{$\modA$ 的投射对象}

讨论作为 $\modA$ 的投射模的基本者。对各自的 $e_i\in A\simeq\Hom_{\modA}(A,A)$，
\[
P(i):=e_i\cdot A\in\modA
\]
是自由 $A$ 模 $A$ 的直和成分，故是投射的。它们之间的射的空间为
\[
\Hom_{\modA}(P(i),P(j))=e_j\cdot\Hom_{\modA}(A,A)\cdot e_i\simeq e_jAe_i,
\]
另一方面，由 4.8.3 节中 $\Ext_{\modA}$ 的性质，当 $k\neq 0$ 时
\begin{equation}
\dim_K\bigl(\Hom_{\Db(\modA)}(P(i),T^{k}(P(j)))\bigr)=0 \tag{6.4}
\end{equation}
成立。也就是说，$(P(1),\dots,P(n))$ 是 $\Db(\modA)$ 的强例外对象系。实际上可知 $(P(1),\dots,P(n))$ 是强例外生成系（系 6.40）。试就各 $P(k)$ 来考虑 $M=\bigoplus_{i\in\Delta_0}M_i$ 的分解。
\begin{equation}
P(k)\cdot e_j=e_k\cdot A\cdot e_j=\Hom_{\modA}(P(j),P(k)) \tag{6.5}
\end{equation}
由此可知：当 $j>k$ 时 $P(k)\cdot e_j=0$；当 $j=k$ 时 $P(k)\cdot e_k=\langle\rho\mid s(\rho)=t(\rho)=k\rangle=\langle e_k\rangle\simeq S(k)$。作为对应于 (6.3) 的短正合列，得到
\begin{equation}
0\to F^{k-1}(P(k))\to P(k)\to S(k)\to 0. \tag{6.6}
\end{equation}

% [Tutor Note: （1）底本 p.158 第 5 行将上式印作 $\Hom_{\modA}(P(i),P(j))=e_j\cdot\Hom_{\modA}(A,A)\cdot e_j\simeq e_jAe_i$；据作者正誤表（2014-05-20 版）订正为 $\cdots\cdot e_i\simeq e_jAe_i$，此处按订正后的内容翻译（否则左端与 $e_jAe_i$ 不一致；由 $\Hom_{\modA}(e_iA,e_jA)$ 中「先右乘 $e_i$、再左乘 $e_j$」亦应为 $e_j\Hom(A,A)e_i$）。（2）(6.5) 之后底本印作「$j<k$ では $P(k)\cdot e_j=0$」；据同一正誤表订正为「$j>k$ では $P(k)\cdot e_j=0$」，此处按订正后的内容翻译。验证：$e_kAe_j$ 非零当且仅当存在从 $k$ 到 $j$ 的路，而箭都满足 $t(\alpha)<s(\alpha)$，故路只能使标号下降，于是 $k\to j$ 的路存在蕴含 $k>j$；故 $j\ge k$ 时 $e_kAe_j=0$，亦即 $j>k$ 时为零（$j=k$ 时按文中另论）。]

现在，令 $\langle P(1),\dots,P(n)\rangle\subset\Db(\modA)$ 为包含 $P(1),\dots,P(n)$ 的 $\Db(\modA)$ 的最小全三角子范畴（定义 5.19）。有限生成右 $A$ 模 $M\in\modA$ 在右 $A$ 模的范畴 $\mathrm{Mod}\text{-}A$ 中具有投射分解（命题 3.36），而 $M$ 具有从有限生成自由 $A$ 模 $F$（即投射模）出发的满射 $\pi\colon F\to M$；由于 $\modA$ 是阿贝尔范畴，故 $M$ 的投射分解可以在 $\modA$ 中取到。进一步，以下引理 6.38 意味着其投射分解的长度是有限的。

\begin{definition}{6.37}
固定 $j=1,\dots,n$。有限生成右 $A$ 模 $M\in\modA$ 具有长度为 $j$ 的投射分解
\begin{equation}
0\leftarrow M\leftarrow P^{j}\leftarrow P^{j-1}\leftarrow\cdots\leftarrow P^{1}\leftarrow 0 \tag{6.7}
\end{equation}
且各 $P^{i}$ 都是 $P(1),P(2),\dots,P(j)$ 的有限直和或零对象时，记作 $M\in\cP^{\le j}$。
\end{definition}

\begin{lemma}{6.38}
对任意有限生成右 $A$ 模 $M\in\modA$ 与每个 $j=1,\dots,n$，$F^{j}(M)\in\cP^{\le j}$ 成立。
\end{lemma}

% [原文 p.159]
\begin{proof}
用归纳法证明。$j=1$ 的情形题意已满足。事实上，$F^{1}(M)$ 是 $S(1)=P(1)$ 的直和，故存在投射分解 $0\leftarrow F^{1}(M)\leftarrow P(1)^{\oplus\bullet}\leftarrow 0$。设 $j=k-1$ 时 $F^{k-1}(M)\in\cP^{\le(k-1)}$ 成立。此时考虑 $F^{k}(M)$ 的投射分解。由短正合列 (6.3) 与命题 5.33，在 $\Db(\modA)$ 中 $F^{k}(M)\simeq C\bigl(T^{-1}(M_k)\to F^{k-1}(M)\bigr)$，而 $M_k$ 与 $S(k)$ 的直和同型。这里由假设有 $F^{k-1}(M)\in\cP^{\le(k-1)}$。由它和短正合列 (6.6) 有 $S(k)\in\cP^{\le k}$，即 $M_k\in\cP^{\le k}$。由这些事实，$F^{k}(M)$ 是把 $F^{k-1}(M),M_k$ 替换为对应的投射对象所成复形在 $K^{b}(\cP)$ 中的映射锥而得到的，故得 $F^{k}(M)\in\cP^{\le k}$。
\end{proof}

把此引理用于 $M=F^{n}(M)$，$M\in\modA$ 的情形，可知其投射维数 $\mathrm{pd}(M)\le n$；由该事实与定理 5.46 得到以下。

\begin{theorem}{6.39}
$A$ 的同调维数不超过 $n$，且关于三角范畴的嵌入列
\[
\langle P(1),\dots,P(n)\rangle\subset K^{b}(\cP)\subset\Db(\modA)
\]
三个三角范畴三角等价。
\end{theorem}

把它与 (6.4) 合起来便得到以下结论。

\begin{corollary}{6.40}
$(P(1),\dots,P(n))$ 构成 $\Db(\modA)$ 的强例外生成系。
\end{corollary}

又，由引理 6.38，对任意投射对象 $P$ 都有 $P\in\cP^{\le n}$，故把对应的投射分解写作
\[
0\leftarrow P\leftarrow P^{n}\leftarrow P^{n-1}\leftarrow\cdots\leftarrow P^{1}\leftarrow 0
\]
时，由 $P$ 的投射性，$\id_P\colon P\to P$ 经过 $P^{n}$ 分解。这意味着 $P$ 与 $P^{n}$ 的直和成分同型。也就是说得到以下。

\begin{corollary}{6.41}
$\modA$ 的不可分解投射对象，除同型外仅限于 $P(1),\dots,P(n)$ 之一。
\end{corollary}

另一方面，把由 $P(1),\dots,P(n)$ 构成的 $\modA$ 的全子范畴视为微分为平凡的 DG 范畴。记之为 $\cA$，考虑三角 DG 范畴 $Tw(\cA)$、三角范畴 $Tr(\cA)$。

\begin{corollary}{6.42}
存在三角范畴等价 $Tr(\cA)\simeq\Db(\modA)$。
\end{corollary}

\begin{proof}
对 $K^{b}(\cP)$ 的对象 $\{P^{i},d^{i}\}_{i\in\Z}$，把 $\cX:=\bigoplus_{i\in\Z}P^{i}$ 表示为 $P(1),\dots,P(n)$ 与其次数平移的直和，并由各 $d^{i}$ 构造对应的矩阵 $\Phi\in\Hom^{1}_{DG(\cP)}(\cX,\cX)$；则 $(\cX,\Phi)$ 构成扭复形一事，由 $\{P^{i},d^{i}\}_{i\in\Z}$ 是复形而得；$(\cX,\Phi)$ 构成单侧扭复形一事，由 $(P(1),\dots,P(n))$ 构成强例外对象系而得。（也可视为定理 5.56 的推论。）
\end{proof}

\subsection{$\modA$ 的内射对象}

设 $A\text{-}\mathrm{mod}$ 为有限生成左 $A$ 模所成阿贝尔范畴。两个反变函子

% [原文 p.160]
$\delta\colon\modA\to A\text{-}\mathrm{mod}$、$D\colon A\text{-}\mathrm{mod}\to\modA$ 由
\[
\delta:=\Hom_{\modA}(\ ,A),\qquad D:=\Hom_K(\ ,K)
\]
定义。它们构成到 $\mathrm{vect}(K)$ 的反变函子一事，已在更一般的状况下于 4.1.2 节讨论过。不过，对 $\delta$ 要把 $A\in\modA$、对 $D$ 要把 $\modA$ 的对象或 $K$ 视为 $\mathrm{vect}(K)$ 的对象。进一步，$M\in\modA$ 所具有的 $A$ 的右作用定义了 $\delta(M)$ 上的 $A$ 的左作用，$N\in A\text{-}\mathrm{mod}$ 所具有的 $A$ 的左作用定义了 $D(N)$ 上的 $A$ 的右作用，并可确认射与这些作用相容。

同样地，由
\[
D^{*}:=\Hom_K(\ ,K),\qquad \delta^{*}:=\Hom_{\modA}(A,\ )
\]
定义反变函子 $D^{*}\colon\modA\to A\text{-}\mathrm{mod}$、$\delta^{*}\colon A\text{-}\mathrm{mod}\to\modA$。关于 $D,D^{*}$，注意由向量空间的对偶性有
\[
D^{*}\circ D\overset{\mathrm{n.e.}}{\simeq}\id_{A\text{-}\mathrm{mod}},\qquad
D\circ D^{*}\overset{\mathrm{n.e.}}{\simeq}\id_{\modA}
\]
成立（参看 [2, p12]）。

\begin{definition}{6.43}
称 $D\circ\delta\colon\modA\to\modA$ 为中山函子（Nakayama functor）。
\end{definition}

\begin{theorem}{6.44}
$\modA$ 的不可分解内射模与
\[
I(i):=D\circ\delta(P(i)),\qquad i=1,\dots,n
\]
之一同型。
\end{theorem}

\begin{proof}
由 $D$ 的对偶性（射的方向反转），$M\in\modA$ 是内射模一事与 $D(M)\in A\text{-}\mathrm{mod}$ 是投射模一事等价。另一方面，由
\[
\check P(i):=\delta(P(i))=\Hom_{\modA}(e_iA,A)\simeq A\cdot e_i\in A\text{-}\mathrm{mod}
\]
知 $\check P(1),\dots,\check P(n)$ 给出 $A\text{-}\mathrm{mod}$ 的不可分解投射模的列表。因此，$D(\check P(i))=I(i)$ 是不可分解内射模的列表便可得到。
\end{proof}

\subsection{单对象与投射对象、内射对象的关系}

以下继续设 $A$ 为某有限型有序箭图 $\vec{\Delta}$ 的带关系式路代数，并设 $\sharp(\vec{\Delta}_0)=n$。关于 $\Db(\modA)$ 的强例外生成系 $(P(1),\dots,P(n))$ 上的变异所给出的不可分解投射对象与单对象之间，有以下关系。

\begin{proposition}{6.45}
关于 $\Db(\modA)$ 的强例外生成系 $(P(1),\dots,P(n))$，有
\[
\mathcal L(P(1),\dots,P(n))=(S(n),\dots,S(1)),
\]
\[
\mathcal L:=L'_2(L'_2L'_3)(L'_2L'_3L'_4)\cdots(L'_2\cdots L'_n)
\]
成立。这里设 $L'_i:=TL_i$（$T\colon\cT\to\cT$ 是平移函子）。
\end{proposition}

更具体地说，即
\[
(S(n),\dots,S(3),S(2),S(1))=(T^{n-1}L^{n-1}(P(n)),\dots,T^{2}L^{2}(P(3)),TL(P(2)),P(1)).
\]

% [原文 p.161]
\begin{proof}
对各单对象 $S(k)$，$k=1,\dots,n$，由 $S(k)=F^{k}(S(k))$ 与引理 6.38 知 $S(k)\in\cP^{\le k}$。具体构造对应的投射分解，并验证它与 $T^{k}L^{k}(P(k))$ 一致。把投射分解表为
\begin{equation}
0\to P^{1}\to\cdots\to P^{k-3}\xrightarrow{\tilde\gamma}P^{k-2}\xrightarrow{\tilde\beta}P^{k-1}\xrightarrow{\tilde\alpha}P^{k}\to S(k)\to 0, \tag{6.8}
\end{equation}
并构造 $P^{k},P^{k-1},\dots,\tilde\alpha,\tilde\beta,\dots$。由短正合列 (6.6)，首先取
\[
P^{k}=P(k),
\]
并考虑 $F^{k-1}(P(k))\in\cP^{\le(k-1)}$ 的投射分解。由 (6.5) 有 $F^{k-1}(P(k))\simeq\bigoplus_{i\le k-1}e_k\cdot A\cdot e_i$。注意以 $k$ 为起点、以 $i$ 为终点的路必定可写成以 $k$ 为起点的长度 $1$ 的路 $\alpha$ 与其他路的合成；令 $P^{k-1}:=\bigoplus_{s(\alpha)=k}P(t(\alpha))$、$\tilde\alpha:=\bigoplus_{s(\alpha)=k}\alpha$，则得到满射
\[
P^{k-1}\xrightarrow{\tilde\alpha}F^{k-1}(P(k))\to 0.
\]
进一步考虑此 $\Ker(\tilde\alpha)$。它与由以 $k$ 为起点的关系式 $\rho_1,\dots,\rho_l$ 张成的向量空间同型。对各 $\rho_j$，取对应于其终点 $t(\rho_j)$ 的 $P(t(\rho_j))$，令 $P^{k-2}:=\bigoplus_{j=1}^{l}P(t(\rho_j))$，则得到满射 $P^{k-2}\to\Ker(\tilde\alpha)$。令 $\tilde\beta\colon P^{k-2}\to P^{k-1}$ 为对应的映射，如此便得到正合列
\[
P^{k-2}\xrightarrow{\tilde\beta}P^{k-1}\xrightarrow{\tilde\alpha}P(k)\to S(k)\to 0.
\]
继续这样的作业便得到投射分解 (6.8)。汇总而言，取
\begin{itemize}[leftmargin=1.6em,itemsep=1pt,topsep=2pt]
\item $P^{k-1}:=\bigoplus_{s(\alpha)=k}P(t(\alpha))$，$\tilde\alpha:=\bigoplus_{s(\alpha)=k}\alpha$，
\item $P^{k-2}:=\bigoplus_{s(\rho)=k}P(t(\rho))$，$\tilde\alpha\circ\tilde\beta=\bigoplus_{s(\rho)=k}\rho$，
\item $P^{k-3}:=\bigoplus_{s(\alpha)=k}\bigl(\bigoplus_{s(\rho)=t(\alpha)}P(t(\rho))\bigr)$，$\tilde\beta\circ\tilde\gamma=\bigoplus_{s(\alpha)=k}\bigl(\bigoplus_{s(\rho)=t(\alpha)}\rho\bigr)$，
\item $\cdots$
\end{itemize}
即可。这里，$\rho$ 是定义 $A$ 的 $K\vec{\Delta}$ 的关系式，而 $\tilde\alpha\circ\tilde\beta,\tilde\beta\circ\tilde\gamma,\dots$ 视为 $\mathrm{Mat}(K\vec{\Delta})$ 的元。

% [Tutor Note: 底本 p.161 第 $-14$、$-13$ 行将上列第 2、3 条分别印作 $\tilde\beta\circ\tilde\alpha=\bigoplus_{s(\rho)=k}\rho$、$\tilde\gamma\circ\tilde\beta=\bigoplus_{s(\alpha)=k}(\bigoplus_{s(\rho)=t(\alpha)}\rho)$；据作者正誤表（2014-05-20 版）订正为 $\tilde\alpha\circ\tilde\beta$、$\tilde\beta\circ\tilde\gamma$，此处按订正后的内容翻译。理由：$\tilde\alpha:P^{k-1}\to P^k$、$\tilde\beta:P^{k-2}\to P^{k-1}$，故从 $P^{k-2}$ 到 $P^k$ 的合成应为 $\tilde\alpha\circ\tilde\beta$；同理 $\tilde\beta\circ\tilde\gamma$。另，底本第 $-11$ 行（汇总句）仍印作「$\tilde\beta\circ\tilde\alpha,\tilde\gamma\circ\tilde\beta,\dots$」，正誤表未列该行，此处照原印译出；读者应依上述订正理解为 $\tilde\alpha\circ\tilde\beta,\tilde\beta\circ\tilde\gamma$。]

其次，构造 $T^{k-1}L^{k-1}(P(k))$。这可通过在 $K^{b}(\cP)$ 中反复作左变异而得到，但由系 6.42，用 $Tw(\cA)$ 中的单侧扭复形来表示。不过要用下三角矩阵表示（参看注意 5.51）。令 $\dim_K\bigl(\Hom_{\modA}(P(k-1),P(k))\bigr)=d_{k-1}$，把对应的映射记为 $\alpha\colon(P(k-1))^{\oplus d_{k-1}}\to P(k)$，则
\[
TL(P(k))=\bigl((TP(k-1))^{\oplus d_{k-1}}\oplus P(k),\ \Phi_{TL(P(k))}\bigr),
\]
\[
\Phi_{TL(P(k))}:=
\begin{pmatrix}
0 & 0\\
\alpha & 0
\end{pmatrix}
\]
成为。进一步，对定义 $TL(P(k))$ 的正合三角序列施加上同调函子 $\Hom_{\Db(\modA)}(P(k-2),\ )$，便得到

% [原文 p.162]
\[
\Hom_{\Db(\modA)}(P(k-2),TL(P(k)))\simeq\langle\alpha\in\vec{\Delta}_1\mid s(\alpha)=k,\ t(\alpha)=k-2\rangle,
\]
\[
\Hom_{\Db(\modA)}(TP(k-2),TL(P(k)))\simeq\Ker\bigl(\Hom_{\modA}(P(k-2),\alpha)\bigr),
\]
\[
\Hom_{\Db(\modA)}(T^{k}P(k-2),TL(P(k)))=0,\qquad k\neq 0,1,
\]
并由此得到：把 $TP(k-2)$ 的 $\dim_K(\langle\alpha\in\vec{\Delta}_1\mid s(\alpha)=k,\ t(\alpha)=k-2\rangle)$ 个直和出发的次数 $1$ 的射记为 $\alpha\colon T(P(k-2))^{\oplus\bullet}\to P(k)$，把 $T^{2}P(k-2)$ 的 $\dim_K\bigl(\Ker(\Hom_{\modA}(P(k-2),\alpha))\bigr)$ 个直和出发的次数 $1$ 的射记为 $\beta\colon T^{2}(P(k-2))^{\oplus\bullet}\to T(P(k-1))^{\oplus\bullet}$，则 $T^{2}L^{2}(P(k))$ 为
\[
\bigl((T^{2}(P(k-2))^{\oplus\bullet}\oplus T(P(k-2))^{\oplus\bullet})\oplus T(P(k-1))^{\oplus\bullet}\oplus P(k),\ \Phi_{T^{2}L^{2}(P(k))}\bigr),
\]
\[
\Phi_{T^{2}L^{2}(P(k))}:=
\left(\begin{array}{cc|c|c}
0 & 0 & 0 & 0\\
0 & 0 & 0 & 0\\ \hline
\beta & 0 & 0 & 0\\ \hline
0 & \alpha & \alpha & 0
\end{array}\right)
\]
成为。同样地，$T^{k-1}L^{k-1}(P(k))=:(\cX,\Phi)$ 为
\[
\cX:=\bigl(T^{k-1}(P(1))^{\oplus\bullet}\oplus\cdots\oplus T(P(1))^{\oplus\bullet}\bigr)\oplus\cdots
\]
\[
\qquad\cdots\oplus\bigl(T^{2}(P(k-2))^{\oplus\bullet}\oplus T(P(k-2))^{\oplus\bullet}\bigr)\oplus T(P(k-1))^{\oplus\bullet}\oplus P(k),
\]
\[
\Phi:=
\left(\begin{array}{c|ccc|cc|c|c}
\ddots & \vdots & \vdots & \vdots & \vdots & \vdots & \vdots & \vdots\\ \hline
\cdots & 0 & 0 & 0 & 0 & 0 & 0 & 0\\
\cdots & 0 & 0 & 0 & 0 & 0 & 0 & 0\\
\cdots & 0 & 0 & 0 & 0 & 0 & 0 & 0\\ \hline
\cdots & \gamma & 0 & 0 & 0 & 0 & 0 & 0\\
\cdots & 0 & \beta & 0 & 0 & 0 & 0 & 0\\ \hline
\cdots & 0 & \beta & 0 & \beta & 0 & 0 & 0\\ \hline
\cdots & 0 & 0 & \alpha & 0 & \alpha & \alpha & 0
\end{array}\right)
\]
成为；而关于这些 $T^{\bullet}P(\bullet)$，$T$ 的次数为 $i$ 者们的直和与复形 (6.8) 中的 $P^{k-i}$ 一致，由上面的 $\Phi$ 把 $\alpha$ 们汇总为 $\tilde\alpha$、$\beta$ 们汇总为 $\tilde\beta$ 等，便得到复形 (6.8)（详细略）。
\end{proof}

\begin{corollary}{6.46}
单对象 $(S(n),\dots,S(2),S(1))$ 构成 $\Db(\modA)$ 的例外生成系。
\end{corollary}

\begin{proof}
由 $(S(n),\dots,S(2),S(1))$ 是由例外对象系 $(P(1),\dots,P(n))$ 变异而得到，即得。
\end{proof}

\begin{corollary}{6.47}
$\Ext^{1}_{\modA}(S(i),S(j))\simeq\langle\alpha\in\Delta_1\mid s(\alpha)=i,\ t(\alpha)=j\rangle$ 成立。
\end{corollary}

\begin{proof}
用投射分解 (6.8) 直接计算 $\Ext^{1}_{\modA}$ 即可。
\end{proof}

\begin{corollary}{6.48}
若存在有限型有序箭图 $\vec{\Delta}$ 使 $A=K\vec{\Delta}$，则 $A$ 的同调维数为 $1$，且当 $r\neq 1$ 时 $\Ext^{r}(S(i),S(j))=0$。
\end{corollary}

% [原文 p.163]
\begin{proof}
由各 $S(i)$ 具有长度 $1$ 的投射分解，且追溯引理 6.38 的证明可知任意对象 $M\in\modA$ 都具有长度 $1$ 的投射分解。由此便有 $\Ext^{r}(S(i),S(j))=0$（$r\neq 1$）。
\end{proof}

在不可分解内射模与单模之间也存在同样的关系。

\begin{proposition}{6.49}
关于 $\Db(\modA)$ 的例外生成系 $(S(n),\dots,S(1))$，有
\[
\mathcal L(S(n),\dots,S(1))=(I(1),\dots,I(n)),
\]
\[
\mathcal L:=L'_2(L'_2L'_3)(L'_2L'_3L'_4)\cdots(L'_2\cdots L'_n),\qquad L'_i:=TL_i
\]
成立。
\end{proposition}

\begin{proof}
对有限生成左 $A$ 模的范畴 $A\text{-}\mathrm{mod}$ 的不可分解投射对象 $\check P(i):=A\cdot e_i$，$i=1,\dots,n$（参看 6.5.3 节），$(\check P(n),\dots,\check P(1))$ 构成 $\Db(A\text{-}\mathrm{mod})$ 的强例外生成系；与 $\Db(\modA)$ 的情形同样地，
\[
\mathcal L(\check P(1),\dots,\check P(n))=(\check S(n),\dots,\check S(1))
\]
成立。令 $R'_i:=T^{-1}R_i$，则因 $R'_i\circ L'_{i+1}=1$，所以令
\[
\mathcal R:=R'_1(R'_2R'_1)(R'_3R'_2R'_1)\cdots(R'_{n-1}\cdots R'_1),\qquad R'_i:=T^{-1}R_i
\]
时 $\mathcal R\circ\mathcal L=1$，并得到
\[
\mathcal R(\check S(n),\dots,\check S(1))=(\check P(1),\dots,\check P(n)).
\]
对这些取双对 $D\colon A\text{-}\mathrm{mod}\to\modA$，则射的方向反转，$R_i$ 与 $L_{i+1}$ 互换，从而得到题意。
\end{proof}

下面用例来确认 $\Db(\modA)$ 与 $\Db(A\text{-}\mathrm{mod})$ 之间的关系。

\begin{example}{6.50}
试考虑例 6.12 的箭图 $\vec{\Delta}$ 的路代数 $A$。作为单对象的投射分解得到以下。
\[
\vec{\Delta}=
\Biggl(
\begin{tikzpicture}[baseline=-0.5ex]
  \node[draw,circle,inner sep=1pt,label=left:{$1$}] (v1) at (0,0) {};
  \node[draw,circle,inner sep=1pt,label=right:{$2$}] (v2) at (2.2,0) {};
  \node[draw,circle,inner sep=1pt,label=right:{$3$}] (v3) at (3.1,1.5) {};
  \node[draw,circle,inner sep=1pt,label=right:{$4$}] (v4) at (3.1,-1.5) {};
  \draw[-{Stealth}] (v2) -- node[below]{$\alpha_{21}$} (v1);
  \draw[-{Stealth}] (v3) -- node[above left]{$\alpha_{32}$} (v2);
  \draw[-{Stealth}] (v4) -- node[below left]{$\alpha_{42}$} (v2);
\end{tikzpicture}
\Biggr),
\]
\[
S(1)=P(1),\qquad
0\leftarrow S(2)\leftarrow P(2)\leftarrow P(1)\leftarrow 0,
\]
\[
0\leftarrow S(3)\leftarrow P(3)\leftarrow P(2)\leftarrow 0,\qquad
0\leftarrow S(4)\leftarrow P(4)\leftarrow P(2)\leftarrow 0.
\]
另一方面，考虑 $A=K\vec{\Delta}$ 的有限生成左模的范畴 $A\text{-}\mathrm{mod}$，则其单对象 $\check S(i)$ 的、由不可分解投射左模 $\check P(1),\check P(2),\check P(3),\check P(4)$ 给出的投射分解为
\[
0\leftarrow\check S(1)\leftarrow\check P(1)\leftarrow\check P(2)\leftarrow 0,
\]
\[
0\leftarrow\check S(2)\leftarrow\check P(2)\leftarrow(\check P(3)\oplus\check P(4))\leftarrow 0,
\]
\[
\check S(3)=\check P(3),\qquad \check S(4)=\check P(4).
\]
取此双对 $D$，便得到 $S(i)\in\modA$ 的内射分解。

% [原文 p.164]
\[
0\to S(1)\to I(1)\to I(2)\to 0,
\]
\[
0\to S(2)\to I(2)\to(I(3)\oplus I(4))\to 0,
\]
\[
S(3)=I(3),\qquad S(4)=I(4).
\]
由这些单对象的投射分解、内射分解，也可以读出对应的变异所给出的关系。
\end{example}

\subsection{变异与 Serre 函子}

$A$ 仍如以前，为某有限型有序箭图的带关系式 $K$ 上路代数。$\Db(\modA)$ 具有强例外生成系 $(P(1),\dots,P(n))$，故由定理 5.74，$\Db(\modA)$ 具有 Serre 函子（定义 5.71）。另一方面，关于这个有限维代数的导出范畴，可以具体给出 Serre 函子。中山函子（定义 6.43）自然地提升为对应的导出范畴上的函子。把它们用同一记号 $\delta\colon\Db(\modA)\to\Db(A\text{-}\mathrm{mod})$、$D\colon\Db(A\text{-}\mathrm{mod})\to\Db(\modA)$、$D\circ\delta\colon\Db(\modA)\to\Db(\modA)$ 表示。$D\circ\delta$ 构成自等价函子。事实上，考虑 $\delta^{*}\circ D^{*}\colon\Db(\modA)\to\Db(\modA)$ 时，它们由定义是三角函子，而
\[
(\delta^{*}\circ D^{*})\circ(D\circ\delta)\overset{\mathrm{n.e.}}{\simeq}\id_{\Db(\modA)}\overset{\mathrm{n.e.}}{\simeq}(D\circ\delta)\circ(\delta^{*}\circ D^{*})
\]
一事只要在各 $P(i)$ 上确认即可。

\begin{theorem}{6.51}
中山函子 $\nu:=D\circ\delta\colon\Db(\modA)\to\Db(\modA)$ 构成 $\Db(\modA)$ 的 Serre 函子。
\end{theorem}

由 Serre 函子的唯一性，若 $D\circ\delta$ 是 Serre 函子，则被保障为与定理 5.74 所得者自然同值。也就是说，应有
\begin{equation}
D\circ\delta(P(i))=I(i)\simeq T^{n-1}L^{n-1}(P(i)) \tag{6.9}
\end{equation}
而这是实际成立的。首先由前节的命题 6.45、6.49 有
\[
(\mathcal L)^{2}(P(1),\dots,P(n))=(I(1),\dots,I(n)),
\]
故 (6.9) 与等式
\[
(\mathcal L)^{2}(P(1),\dots,P(n))=(T^{n-1}L^{n-1}(P(1)),\dots,T^{n-1}L^{n-1}(P(n)))
\]
等价；而这正是把由变异性质（命题 5.69）得到的例外对象系 $(E_1,\dots,E_n)$ 上的等式
\[
(\mathcal L)^{2}=(L'_2\cdots L'_n)^{n},\qquad L'_i:=TL_i
\]
作用于 $(P(1),\dots,P(n))$ 者本身。

由这一事实 (6.9) 与以下命题，定理 6.51 便得到。

\begin{proposition}{6.52}
对具有例外生成系 $(E_1,\dots,E_n)$ 的 $K$ 上有限型三角范畴 $\cT$，若与平移函子可换的自等价函子 $\nu\colon\cT\to\cT$ 对每个 $i=1,\dots,n$ 满足 $\nu(E_i)\simeq T^{n-1}L^{n-1}(E_i)$，则 $\nu$ 是 Serre 函子。
\end{proposition}

% [原文 p.165]
\begin{proof}
设 $\cT$ 的正合三角序列
\begin{equation}
A\to B\to C\to A[1] \tag{6.10}
\end{equation}
与任意 $i\in\Z$ 之间存在同型
\[
\Hom_{\cT}(A[i],X)\simeq\Hom_{\cT}(X,\nu(A[i]))^{*},\qquad
\Hom_{\cT}(C[i],X)\simeq\Hom_{\cT}(X,\nu(C[i]))^{*},
\]
且这些定出自然同值，此时证明关于 $B$ 的同样自然同值存在。由于 $\nu$ 是自等价函子，故 $\mu:=\nu^{-1}$ 存在；对任意对象 $X,Y\in\cT$ 有 $\Hom_{\cT}(X,\nu(Y))^{*}\simeq\Hom_{\cT}(\mu(X),Y)^{*}$，故上述假定与以下条件等价：
\[
\Hom_{\cT}(A[i],X)\simeq\Hom_{\cT}(\mu(X),A[i])^{*},\qquad
\Hom_{\cT}(C[i],X)\simeq\Hom_{\cT}(\mu(X),C[i])^{*}.
\]
对正合三角序列 (6.10) 施加反变上同调函子 $\Hom_{\cT}(\ ,X)$、$\Hom_{\cT}(\mu(X),\ )^{*}$，便得到向量空间的图表
\[\begin{tikzcd}[row sep=3.0em, column sep=2.0em]
  \cdots \arrow[r] & \Hom_{\cT}(A,X) \arrow[r] \arrow[d,"\sim"'] & \Hom_{\cT}(B,X) \arrow[r] \arrow[d,dashed,"?"] & \Hom_{\cT}(C,X) \arrow[r] \arrow[d,"\sim"'] & \cdots \\
  \cdots \arrow[r] & \Hom_{\cT}(\mu(X),A)^{*} \arrow[r] & \Hom_{\cT}(\mu(X),B)^{*} \arrow[r] & \Hom_{\cT}(\mu(X),C)^{*} \arrow[r] & \cdots
\end{tikzcd}\]
并可构成使其交换的、点线部分的同型射。它定出自然同值 $\Hom_{\cT}(B,\ )\overset{\mathrm{n.e.}}{\simeq}\Hom_{\cT}(\mu(\ ),B)^{*}\overset{\mathrm{n.e.}}{\simeq}\Hom_{\cT}(\ ,\nu(B))^{*}$，归纳地对任意对象 $B\in\langle E_1,\dots,E_n\rangle\simeq\cT$ 证明了这一自然同值的存在。
\end{proof}

\section{三角范畴的强例外生成系与有限维代数}

以下定理意味着有限型有序箭图的带关系式路代数 $A$ 的导出范畴 $\Db(\modA)$ 出现在各种各样之处。

\begin{theorem}{6.53}
（[7]）当具有 DG 增广的 $K$ 上三角范畴 $\cT$ 具有强例外生成系 $(E_1,\dots,E_n)$ 时，对 $K$ 自同态代数
\[
A:=\Hom_{\cT}(E,E),\qquad E:=\bigoplus_{i=1}^{n}E_i
\]
存在三角范畴等价 $\cT\simeq\Db(\modA)$。
\end{theorem}

由于 $\Db(\modA)$ 是仅由代数 $A$ 定义者，故三角范畴 $\cT$ 仅由代数 $A$ 的信息决定。在这个意义上，定理 6.53 也可以说是定理 5.56 在特别状况下的精密化。

% [原文 p.166]
以下，作为定理 6.53 的一个证明，说明应用下一章所述定理 7.51、特别是系 7.53 者。当具有强例外生成系 $(E_1,\dots,E_n)$ 的 $K$ 上三角范畴 $\cT$ 具有 DG 增广时，设对象 $E_1,\dots,E_n$ 上的 DG 范畴为 $\cC$，其零次上同调所得的范畴为 $H^{0}(\cC)$，则由系 7.53 得到以下事实。
\[
\cT=Tr(\cC)\simeq Tr(H^{0}(\cC))=Tr(\cH),\qquad \cH:=H^{0}(\cC).
\]
这里 $Tr(\cH)$ 是把 $\cH=H^{0}(\cC)$ 视为微分为平凡的 DG 范畴而得到的三角范畴。另一方面，对由 $P(1),\dots,P(n)$ 构成的 $\modA$ 的全子范畴 $\cA$，存在三角范畴等价 $Tr(\cA)\simeq\Db(\modA)$（系 6.42），故由这些事实，以下构成三角范畴等价函子
\begin{equation}
F\colon Tr(\cH)\to Tr(\cA)\simeq\Db(\modA). \tag{6.11}
\end{equation}
现在对象 $X\in Tr(\cH)$ 是由 $E_1,\dots,E_n$ 与其次数平移构成的单侧扭复形 $(X,\Phi)$。对这些，令 $E:=\bigoplus_{i=1}^{n}E_i$，由
\[
F:=\Hom_{Tr(\cH)}(E,\ )
\]
给出。由定义它给出三角函子。特别地，对每个 $i$，因 $\Hom_{\cH}(E,E_i)\simeq P(i)$，故得 $F(E_i)=P(i)$。也就是说，它把 $Tr(\cH)$ 的生成系映到 $Tr(\cA)$ 的生成系。由此证明了 $F$ 给出三角范畴等价。\hfill$\square$

% [Tutor Note: 原书 p.166 此处的证明只以「由定义它给出三角函子」一句带过 $F=\Hom_{Tr(\cH)}(E,\ )$ 的三角性：$Tr(\cH)$ 的三角结构由单侧扭复形给出，而 $\Hom(E,\ )$ 把 $E_i$ 送到 $P(i)$ 并把扭复形 $(X,\Phi)$ 送到 $D^b(\mathrm{mod}\text{-}A)$ 中的「以 $P(i)$ 的平移为分量的复形」，故与平移、映射锥都相容；这正是定理 5.56 的框架。又，文中用到 $\Hom_{\cH}(E,E_i)\simeq P(i)$，即 $\cH=H^0(\cC)$ 上 $E$ 到 $E_i$ 的态射空间同构于 $P(i)=e_iA$，这是由 $(E_1,\dots,E_n)$ 为强例外生成系（$\Hom(E_i,E_j)$ 的集中性）得到的。]

本题在 $A_\infty$ 范畴中的推广将在 7.4.2 节（系 7.68）中讨论。

\section{范畴的图表示}

\subsection{Krull--Schmidt 范畴的 AR 箭图}

注意对 $K$ 上加法范畴 $\cC$ 的每个对象 $X\in\cC$，$\Hom_{\cC}(X,X)$ 构成 $K$ 代数。称 $e\in\Hom_{\cC}(X,X)=:\End_{\cC}(X)$ 是幂等元，是指 $e$ 作为 $K$ 代数是幂等元（定义 6.3）。

\begin{definition}{6.54}
当 $K$ 上加法范畴 $\cC$ 满足以下性质时，称 $\cC$ 具有 Krull--Schmidt 性。
\begin{itemize}[leftmargin=2em,itemsep=1pt,topsep=2pt]
\item[(a)] 对任意对象 $X,Y\in\cC$，$\Hom_{\cC}(X,Y)$ 是 $K$ 上有限维向量空间。
\item[(b)] 对任意对象 $X\in\cC$ 与幂等元 $e\in\Hom_{\cC}(X,X)$，存在某对象 $Y\in\cC$ 与射 $\pi\in\Hom_{\cC}(X,Y)$、$\iota\in\Hom_{\cC}(Y,X)$，满足
\[
e=\iota\circ\pi,\qquad \pi\circ\iota=\id_Y.
\]
\end{itemize}
\end{definition}

关于 $K$ 上加法范畴 $\cC$ 的对象 $X\in\cC$ 的条件

(i) $\End_{\cC}(X)$ 构成局部环（定义 2.24）。

% [原文 p.167]
(ii) $\End_{\cC}(X)$ 的幂等元只有零射或恒等射 $\id_X$。
(iii) $X\in\cC$ 是不可分解对象（定义 4.28）。

关于这些，有 (i) 则 (ii)、(ii) 则 (iii) 成立；而当 $\cC$ 具有 Krull--Schmidt 性时 (i),(ii),(iii) 等价（两者都可以直接证明）。也就是说，在具有 Krull--Schmidt 性的 $K$ 上加法范畴 $\cC$ 中，把不可分解对象定义为满足条件 (i) 或 (ii) 者也无妨。

\begin{theorem}{6.55}
（Krull--Schmidt 定理）$K$ 上 Krull--Schmidt 加法范畴 $\cC$ 的任意对象 $X\in\cC$ 都与有限个不可分解对象之直和 $X_1\oplus\cdots\oplus X_n$ 同型；此时不可分解对象 $X_1,\dots,X_n$ 除同型与顺序的交换外唯一确定。
\end{theorem}

\begin{proof}
若 $\cC$ 是 $K$ 上 Krull--Schmidt 范畴，则对 $X\in\cC$，有限维 $K$ 代数 $\Hom_{\cC}(X,X)$ 中的幂等元只有有限个，特别地由 $1_X$ 的本原正交幂等元分解 $1_X=e_1+\cdots+e_n$ 得到直和分解 $X\simeq X_1\oplus\cdots\oplus X_n$。唯一性利用不可分解对象上的自同态环是局部环一事来证明（略：[24, p92, 定理 3-2-10]、[46, p26, 定理 6.1]、[51, Prop.5.4]、[2, I.4.10] 等）。
\end{proof}

\begin{example}{6.56}
若 $\cC$ 是某域 $K$ 上的阿贝尔范畴且射的空间全都是有限维的（也就是说，定义 6.54 的 (a) 条件已被满足），则 $\cC$ 具有 Krull--Schmidt 性。事实上，对幂等元 $e\colon X\to X$，考虑 $Y:=\Img(e)$ 便得到图表
\[\begin{tikzcd}[row sep=2.8em, column sep=3.0em]
  & Y \arrow[ddl,dashed,"\pi"'] \arrow[ddr,"\iota"] & \\
  & & \\
  X \arrow[rr,"e"] \arrow[drr,"\pi'"'] & & X \\
  & & \Coker(e)
\end{tikzcd}\]
（由 $\pi'\circ e=0$，$e\colon X\to X$ 经过 $\iota\colon Y\to X$ 分解。）特别地，$\iota\colon Y\to X$ 是单射，$\pi\colon X\to Y$ 是满射，且 $\iota\circ\pi=e$。由于 $e$ 是幂等元，故 $(\iota\circ\pi)\circ(\iota\circ\pi)=e^{2}=e=\iota\circ\pi$；把它改写便得到
\[
\iota\circ(\pi\circ\iota-\id_Y)\circ\pi=0.
\]
这里 $\iota$ 是单射、$\pi$ 是满射，故这意味 $\pi\circ\iota-\id_Y=0$，即知定义 6.54 的 (b) 条件被满足。也就是说，对有限维代数 $A$，阿贝尔范畴 $\modA$、$A\text{-}\mathrm{mod}$ 都满足 Krull--Schmidt 性。
\end{example}

对阿贝尔范畴 $\cC$，其导出范畴 $D(\cC)$ 一般未必是阿贝尔范畴，但导出范畴、三角范畴等之中也有许多满足 Krull--Schmidt 性者。下一节将证明：由有序箭图的带关系式路代数 $A$ 的阿贝尔范畴 $\modA$ 所得的有界导出范畴 $\Db(\modA)$ 具有 Krull--Schmidt 性。

% [原文 p.168]
\begin{example}{6.57}
作为 $K$ 上加法范畴不具有 Krull--Schmidt 性的例子，例如有由 $K$ 上偶数维向量空间所成的全子范畴 $\mathrm{vect}^{+}(K)\subset\mathrm{vect}(K)$。它不是阿贝尔范畴一事已在例 4.40 看到。

设 $V\in\mathrm{vect}^{+}(K)$ 为 $2$ 维向量空间。取基 $v_1,v_2$，则 $\Hom_{\mathrm{vect}^{+}(K)}(V,V)$ 可表为以 $K$ 为成分的 $2\times 2$ 矩阵：
\[
(v_1\ v_2)\mapsto(v_1\ v_2)\begin{pmatrix}a&b\\ c&d\end{pmatrix}.
\]
特别地，作为 $\Hom_{\mathrm{vect}^{+}(K)}(V,V)$ 的元，考虑 $\begin{pmatrix}a&b\\ c&d\end{pmatrix}=\begin{pmatrix}1&0\\ 0&0\end{pmatrix}$，则对应的射 $e_1\in\Hom_{\mathrm{vect}^{+}(K)}(V,V)$ 是幂等元。在阿贝尔范畴 $\mathrm{vect}(K)$ 中，$\Img(e_1)$ 作为以 $v_1$ 为基的 $1$ 维向量空间存在，但它不是 $\mathrm{vect}^{+}(K)$ 的对象。因此 $\mathrm{vect}^{+}(K)$ 不具有 Krull--Schmidt 性。（实际上，不存在对象 $V'\in\mathrm{vect}^{+}(K)$ 与满足 $\iota\circ\pi=e_1$、$\pi\circ\iota=\id_{V'}$ 的单射 $\iota\in\Hom_{\mathrm{vect}^{+}(K)}(V',V)$、满射 $\pi\in\Hom_{\mathrm{vect}^{+}(K)}(V,V')$ 一事，只要取 $V'$ 的基把 $\iota,\pi$ 用矩阵表示便知。）
\end{example}

其次，对范畴的图表示所需的射的不可约性下定义。以下，设 $\cC$ 为与小范畴等价的 $K$ 上 Krull--Schmidt 加法范畴。对两个对象 $X,Y\in\cC$，把满足下述条件的射 $f\in\Hom_{\cC}(X,Y)$ 全体记为
\[
\rad_{\cC}(X,Y)\ (\subset\Hom_{\cC}(X,Y)),
\]
其中条件是：对任何不可分解对象 $Z\in\cC$ 与射 $g\colon Z\to X$、$h\colon Y\to Z$，$h\circ f\circ g\in\End_{\cC}(Z)$ 都不是同构射。这个部分集合实际上是 $\Hom_{\cC}(X,Y)$ 的部分向量空间（参看注意 6.59）。进一步，把 $\rad_{\cC}(X,Y)$ 中满足「存在某对象 $Z\in\cC$，使得该元可表为 $\rad_{\cC}(X,Z)$ 的元与 $\rad_{\cC}(Z,Y)$ 的元的合成」者全体所成的部分集合记为
\[
\rad^{2}_{\cC}(X,Y)\ (\subset\rad_{\cC}(X,Y)).
\]
它也成为 $\rad_{\cC}(X,Y)$ 的部分向量空间。事实上，若 $f,g\in\rad^{2}_{\cC}(X,Y)$ 且 $f,g$ 分别经过 $Z,Z'$ 分解，则 $f+g$ 经过 $Z\oplus Z'$ 分解，故 $f+g\in\rad^{2}_{\cC}(X,Y)$。

\begin{definition}{6.58}
对两个对象 $X,Y\in\cC$，称 $\rad_{\cC}(X,Y)\setminus\rad^{2}_{\cC}(X,Y)$ 的元为不可约射（irreducible morphism）。此外，定义商向量空间
\[
\mathrm{Irr}_{\cC}(X,Y):=\rad_{\cC}(X,Y)/\rad^{2}_{\cC}(X,Y).
\]
\end{definition}

\begin{remark}{6.59}
对应于不可分解分解 $X=\bigoplus_{i=1}^{m}X_i$、$Y:=\bigoplus_{j=1}^{n}Y_j$，把 $\Hom_{\cC}(X,Y)=\bigoplus_{i,j}\Hom_{\cC}(X_i,Y_j)$ 表示出来时，$\rad_{\cC}(X,Y)$ 与「由各 $\Hom_{\cC}(X_i,Y_j)$ 的不是同构射的元所成」的 $\Hom_{\cC}(X,Y)$ 的部分向量空间一致。这一点可通过注意各 $\End_{\cC}(X_i),\End_{\cC}(Y_j)$ 是局部环来证明。由此可知 $\rad_{\cC}(X,Y)$ 构成向量空间。特别地，若 $X,Y\in\cC$ 是不可分解对象，则 $\rad_{\cC}(X,Y)\subset\Hom_{\cC}(X,Y)$ 是由非同构射所成的部分向量空间；若 $X$ 与 $Y$ 不同构，则 $\rad_{\cC}(X,Y)=\Hom_{\cC}(X,Y)$。
\end{remark}

% [原文 p.169]
\begin{definition}{6.60}
在 Krull--Schmidt 范畴 $\cC$ 中，AR 箭图（AR-quiver）\footnote{AR 是 Auslander--Reiten。}$\Gamma(\cC)$，是指箭图 $\Gamma(\cC):=(\Gamma_0,\Gamma_1)$，其顶点集合 $\Gamma_0$ 取为不可分解对象 $X\in\cC$ 的同型类 $[X]$，射集合 $\Gamma_1$ 则使得从各 $[X]\in\Gamma_0$ 到 $[Y]\in\Gamma_0$ 的箭有 $\dim_K(\mathrm{Irr}_{\cC}(X,Y))$ 条。
\end{definition}

对 Krull--Schmidt 范畴 $\cC$，其 AR 箭图未必总能完全描述 $\cC$ 的结构，但非常有用。

\subsection{带关系式有序箭图的有界导出范畴的 Krull--Schmidt 性}

\begin{proposition}{6.61}
设 $\vec{\Delta}$ 为有限型有序箭图，$A$ 为箭图 $\vec{\Delta}$ 的某带关系式路代数。此时 $\Db(\modA)$ 具有 Krull--Schmidt 性。
\end{proposition}

\begin{proof}
对由 $P(1),\dots,P(n)$ 构成的 $\modA$ 的全子范畴 $\cA$，由存在三角范畴等价 $Tr(\cA)\simeq\Db(\modA)$（系 6.42），考虑其 DG 增广 $Tw(\cA)$。与 5.5.2 节中关于 $Tr(\cC)$ 满足正合三角序列公理 (Tr3) 的讨论同样地，对任意对象 $\cX:=(\cX,\Phi)\in Tw(\cA)$、$\cX=\tilde X_1\oplus\tilde X_2\oplus\cdots\oplus\tilde X_N$，若 $\Phi$ 的矩阵成分含有同构射，则存在别的对象 $\cX':=(\cX',\Phi')$ 使得 $\cX'=\tilde X'_1\oplus\cdots\oplus\tilde X'_{N'}$、$N'<N$。因此今设 $\cX=(\cX,\Phi)\in Tw(\cA)$，并考虑 $\Phi$ 的矩阵成分不含同构射者。设存在幂等元 $e\in\Hom_{Tr(\cA)}(\cX,\cX)$，则存在射 $\tilde e\in\Hom^{0}_{Tw(\cA)}(\cX,\cX)$ 使
\begin{equation}
(\tilde e)^{2}-\tilde e\in d\bigl(\Hom^{-1}_{Tw(\cA)}(\cX,\cX)\bigr) \tag{6.12}
\end{equation}
成立。现在把 $\cX=P^{\bullet}(1)\oplus\cdots\oplus P^{\bullet}(n)$ 分解表示。这里各 $P^{\bullet}(n)$ 是 $P(n)$ 们与其次数平移的直和，或零对象。按照这个分解表示，把一般的射 $\Psi\in\Hom^{0}_{Tw(\cA)}(\cX,\cX)$ 分解为分块矩阵：
\[
\Psi=
\begin{pmatrix}
\psi_{11} & \cdots & \psi_{1n}\\
\vdots & \ddots & \vdots\\
\psi_{n1} & \cdots & \psi_{nn}
\end{pmatrix},
\qquad
\psi_{ji}\in\Hom^{0}_{Tw(\cA)}(P^{\bullet}(i),P^{\bullet}(j)).
\]
由 $(P(1),\dots,P(n))$ 是例外生成系，当 $i<j$ 时 $\psi_{ij}=0$；对角成分 $\psi_{ii}$ 由零射、恒等射 $\id_{P^{\bullet}(i)}$ 们与其次数平移的直和构成。令 $F^{k}(\cX)$ 为满足 $i-j\ge k$ 的成分 $\psi_{ij}$ 的直和所成的空间，则得到部分向量空间的列
\[
\Hom^{0}_{Tw(\cA)}(\cX,\cX)=F^{0}(\cX)\supset F^{1}(\cX)\supset\cdots\supset F^{n-1}(\cX)\supset F^{n}(\cX)=0.
\]
另一方面，令
\[
\pi\colon\Hom^{0}_{Tw(\cA)}(\cX,\cX)\to\pi\bigl(\Hom^{0}_{Tw(\cA)}(\cX,\cX)\bigr)
\]
为向对角成分的射影，并令 $e_0:=\pi(\tilde e)$。对 (6.12) 的两边取由 $\pi$ 的射影

% [原文 p.170]
来考虑，便知 $e_0$ 是幂等元。注意现在有 $\tilde e-e_0\in F^{1}(\cX)$。由条件 (6.12)，可取得适当的对象 $\cX^{(1)}\in Tr(\cA)$ 与同型映射 $\Psi_{\cX\cX^{(1)}}\in\Hom^{0}_{Tw(\cA)}(\cX,\cX^{(1)})$，使得 $\Psi_{\cX\cX^{(1)}}\circ(\tilde e-e_0)\circ(\Psi_{\cX\cX^{(1)}})^{-1}\in F^{2}(\cX^{(1)})$（详细略）。令
\[
\Psi_{\cX\cX^{(1)}}\circ\tilde e\circ(\Psi_{\cX\cX^{(1)}})^{-1}=:\tilde e^{(1)},\qquad
\Psi_{\cX\cX^{(1)}}\circ e_0\circ(\Psi_{\cX\cX^{(1)}})^{-1}:=e_0^{(1)},
\]
反复同样的手续，便得到对象 $\cX^{(n-1)}$ 与满足 $\tilde e^{(n-1)}-e_0^{(n-1)}\in F^{n}(\cX^{(n-1)})=0$ 的幂等元 $\tilde e^{(n-1)}=e_0^{(n-1)}\in\Hom^{0}_{Tw(\cA)}(\cX^{(n-1)},\cX^{(n-1)})$。这里由 $\cX^{(n-1)}\simeq\cX$ 可知 $\tilde e\in\Hom^{0}_{Tw(\cA)}(\cX,\cX)$ 是幂等元。此时，对应的直和因子在 $Tw(\cA)$ 中存在。
\end{proof}

\subsection{AR 箭图的例子}

对有序有限型箭图 $\vec{\Delta}$ 的路代数 $A=K\vec{\Delta}$ 的阿贝尔范畴 $\modA$、导出范畴 $\Db(\modA)$，就 $A$ 简单的情形举几个例子。带关系式路代数的情形也可以同样考虑，但对应的范畴 $\modA,\Db(\modA)$ 一般会变复杂，故此处不考虑。

\begin{example}{6.62}
考虑例 6.23 的箭图
\[
\vec{\Delta}:=\Bigl(\,
\begin{tikzpicture}[baseline=-0.5ex]
  \node[draw,circle,inner sep=1pt,label=below:{$1$}] (v1) at (0,0) {};
  \node[draw,circle,inner sep=1pt,label=below:{$2$}] (v2) at (1.8,0) {};
  \draw[-{Stealth}] (v2) -- node[above]{$\alpha$} (v1);
\end{tikzpicture}
\,\Bigr).
\]
$\modA$ 的不可分解对象（的同型类）为 $P(1)\simeq S(1)$、$P(2)$、$S(2)$ 这 $3$ 个。作为 $S(2)$ 的投射分解，存在短正合列 $0\leftarrow S(2)\leftarrow P(2)\leftarrow P(1)\leftarrow 0$（参看例 3.18！），由此可以计算 $S(2),P(1),P(2)$ 之间的所有射，$\modA$ 的 AR 箭图如下。
\[\begin{tikzcd}[row sep=2.4em, column sep=4.0em]
  & {[P(2)]} \arrow[dl] \arrow[dr] & \\
  {[P(1)]} & & {[S(2)]}
\end{tikzcd}
\qquad\qquad
\begin{array}{l}
{[P(1)]=[S(1)]},\\
{[P(2)]=[I(1)]},\\
{[S(2)]=[I(2)]}.
\end{array}
\]
$\Db(\modA)$ 的 AR 箭图如下。
\[\begin{tikzpicture}[baseline=-0.5ex,x=1.0cm,y=1.0cm]
  \node (dl) at (-0.3,-0.5) {$\cdots$};
  \node (A) at (1,1.1) {$A^{-1}$};
  \node (B) at (2,-0.5) {$B^{-1}$};
  \node (C) at (3,1.1) {$C^{-1}$};
  \node (A0) at (4,-0.5) {$A^{+0}$};
  \node (B0) at (5,1.1) {$B^{+0}$};
  \node (C0) at (6,-0.5) {$C^{+0}$};
  \node (A1) at (7,1.1) {$A^{+1}$};
  \node (B1) at (8,-0.5) {$B^{+1}$};
  \node (C1) at (9,1.1) {$C^{+1}$};
  \node (dr) at (10.3,-0.5) {$\cdots$};
  \draw[-{Stealth}] (dl) -- (A);
  \draw[-{Stealth}] (A) -- (B);
  \draw[-{Stealth}] (B) -- (C);
  \draw[-{Stealth}] (C) -- (A0);
  \draw[-{Stealth}] (A0) -- (B0);
  \draw[-{Stealth}] (B0) -- (C0);
  \draw[-{Stealth}] (C0) -- (A1);
  \draw[-{Stealth}] (A1) -- (B1);
  \draw[-{Stealth}] (B1) -- (C1);
  \draw[-{Stealth}] (C1) -- (dr);
\end{tikzpicture}
\]
\[
A^{r}:=[T^{r}(P(1))],\qquad B^{r}:=[T^{r}(P(2))],\qquad C^{r}:=[T^{r}(S(2))].
\]
这是把 $\modA$ 的 AR 箭图与其次数平移的拷贝连接起来者。一般地，对同调维数为 $1$ 的有限维代数 $A$，会成这样。
\end{example}

% [Tutor Note: 底本 p.170 例 6.62 印作「$P(1)\simeq S(2)$」；据作者正誤表（2014-05-20 版）订正为「$P(1)\simeq S(1)$」，此处按订正后的内容翻译。核对：由例 6.23 有 $P(1)=e_1A\simeq K\leftarrow 0$，而 $S(1)$ 的定义正是 $\bigl(K\leftarrow 0\bigr)$（见 定義 6.35 及 $S(1)=P(1)$ 一行），故应为 $S(1)$；原印 $S(2)$ 与后文「$P(1),P(2),S(2)$ 三个不可分解对象」重数也不合。]

% [Tutor Note: 例 6.62 的 $\Db(\mathrm{mod}\text{-}A)$ AR 箭图是 $A_2$ 型（两个顶点、一条箭）代数的导出 AR 箭图，为周期 $3$ 的链 $\cdots\to A^{-1}\to B^{-1}\to C^{-1}\to A^{+0}\to B^{+0}\to C^{+0}\to A^{+1}\to\cdots$，画成上下两行交错（上行为 $A^{-1},C^{-1},B^{+0},A^{+1},C^{+1}$，下行为 $B^{-1},A^{+0},C^{+0},B^{+1}$）。可验证：每个顶点恰有一条入箭与一条出箭，正是 Auslander--Reiten 箭图在 $D^b$ 中的 $\mathbb Z$-覆盖。]

\begin{example}{6.63}
对箭图
\[
\vec{\Delta}:=\Bigl(\,
\begin{tikzpicture}[baseline=-0.5ex]
  \node[draw,circle,inner sep=1pt,label=below:{$1$}] (v1) at (0,0) {};
  \node[draw,circle,inner sep=1pt,label=below:{$2$}] (v2) at (1.8,0) {};
  \node[draw,circle,inner sep=1pt,label=below:{$3$}] (v3) at (3.6,0) {};
  \draw[-{Stealth}] (v2) -- node[above]{$\alpha_{21}$} (v1);
  \draw[-{Stealth}] (v3) -- node[above]{$\alpha_{32}$} (v2);
\end{tikzpicture}
\,\Bigr),\qquad
\vec{\Delta}':=\Bigl(\,
\begin{tikzpicture}[baseline=-0.5ex]
  \node[draw,circle,inner sep=1pt,label=below:{$2$}] (w2) at (0,0) {};
  \node[draw,circle,inner sep=1pt,label=below:{$1$}] (w1) at (1.9,0) {};
  \node[draw,circle,inner sep=1pt,label=below:{$3$}] (w3) at (3.8,0) {};
  \draw[-{Stealth}] (w2) -- node[above]{$\alpha_{21}$} (w1);
  \draw[-{Stealth}] (w3) -- node[above]{$\alpha_{31}$} (w1);
\end{tikzpicture}
\,\Bigr),
\]
对应的阿贝尔范畴的 AR 箭图 $\Gamma:=\Gamma(\mathrm{mod}\text{-}K\vec{\Delta})$、$\Gamma':=\Gamma(\mathrm{mod}\text{-}K\vec{\Delta}')$ 如下。
\[\begin{tikzpicture}[baseline=-0.5ex,x=1.15cm,y=1.05cm]
  \node (P3) at (1,2.4) {$[P(3)]$};
  \node (P2) at (0,1.2) {$[P(2)]$};
  \node (I2) at (2,1.2) {$[I(2)]$};
  \node (P1) at (-0.9,0) {$[P(1)]$};
  \node (S2) at (1,0) {$[S(2)]$};
  \node (I3) at (2.9,0) {$[I(3)]$};
  \draw[-{Stealth}] (P1) -- (P2);
  \draw[-{Stealth}] (P2) -- (P3);
  \draw[-{Stealth}] (P3) -- (I2);
  \draw[-{Stealth}] (P2) -- (S2);
  \draw[-{Stealth}] (S2) -- (I2);
  \draw[-{Stealth}] (I2) -- (I3);
  \node at (4.2,1.2) {$\Gamma:$};
\end{tikzpicture}
\qquad\qquad
\begin{tikzpicture}[baseline=-0.5ex,x=1.15cm,y=1.05cm]
  \node (P3) at (0.2,2.4) {$[P(3)]$};
  \node (I2) at (2.8,2.4) {$[I(2)]$};
  \node (P1) at (-0.2,1.2) {$[P(1)]$};
  \node (I1) at (2.2,1.2) {$[I(1)]$};
  \node (P2) at (0.4,0) {$[P(2)]$};
  \node (I3) at (2.8,0) {$[I(3)]$};
  \draw[-{Stealth}] (P1) -- (P3);
  \draw[-{Stealth}] (P3) -- (I1);
  \draw[-{Stealth}] (I1) -- (I2);
  \draw[-{Stealth}] (P1) -- (P2);
  \draw[-{Stealth}] (P2) -- (I1);
  \draw[-{Stealth}] (I1) -- (I3);
  \node at (4.2,1.2) {$\Gamma':$};
\end{tikzpicture}
\]

% [原文 p.171]
这些导出范畴是三角等价的。AR 箭图如下：
\[\begin{tikzpicture}[baseline=-0.5ex,x=0.82cm,y=0.9cm]
  \foreach \i in {0,...,6}{\node[circle,fill,inner sep=1.3pt] at (2*\i,2) {};}
  \foreach \i in {0,...,5}{\node[circle,fill,inner sep=1.3pt] at (2*\i+1,1) {};}
  \foreach \i in {0,...,6}{\node[circle,fill,inner sep=1.3pt] at (2*\i,0) {};}
  \foreach \i in {0,...,5}{
    \draw[-{Stealth}] (2*\i,2) -- (2*\i+1,1);
    \draw[-{Stealth}] (2*\i,0) -- (2*\i+1,1);
    \draw[-{Stealth}] (2*\i+1,1) -- (2*\i+2,2);
    \draw[-{Stealth}] (2*\i+1,1) -- (2*\i+2,0);
  }
  \node at (-1.1,1) {$\cdots\cdots$};
  \node at (13.1,1) {$\cdots\cdots$};
\end{tikzpicture}\]
实际上，在这个三角范畴中，作为箭图 $\vec{\Delta}$ 所附带的强例外生成系 $(P(1),P(2),P(3))$ 的变异 $R_1(P(1),P(2),P(3))$，便得到箭图 $\vec{\Delta}'$ 所附带的强例外生成系 $(P(1),P(2),P(3))$。\hfill$\square$
\end{example}

关于 $\modA$ 的 AR 箭图，[2] 中有许多例子。$\Db(\modA)$ 的 AR 箭图的例子可见于 [18, Chap I.5.] 等。

\para{文献。}关于有限维代数表示论的基本全貌，可举岩永·佐藤 [24]、永尾·津岛 [46]、Assem--Simson--Skowro\'nski [2]（有 vol.1 到 vol.3）、Auslander--Reiten--Smal\o{} [5]、Gabriel--Roiter [16]；包含与之相伴的三角范畴结构者，可举 Happel [18]。作为与有限维代数相关的交换环上模的范畴之书，也有 Yoshino [57]。6.5 节的内容主要沿 Bondal 的原始论文 [7] 的脉络。所谓 AR 箭图，多指把通常称为 AR 变换者图示化所得者。这个 AR 变换与本书所讨论的 Serre 函子密切相关，相关的 AR 序列、AR 三角序列的理论是非常重要的东西（[18] 等）。

% [Tutor Note: p.171 的 AR 箭图是 $A_3$ 型（$1\leftarrow 2\leftarrow 3$）代数的 $D^b$ 之 AR 箭图：三行点阵、每行 7 个（中行 6 个），每个中行顶点有 4 条箭（上、下各一入一出的「菱形」）；这是在 $\mathbb Z$-覆盖意义下的周期 4 模式，与原书图像一致。文中「变异 $R_1(P(1),P(2),P(3))$ 给出 $\vec\Delta'$ 的强例外生成系」正是「变异改变箭图的定向」这一现象的最简单实例，可与 6.5.4 节命題 6.45 的 $\mathcal L$ 公式对照验证。]